Kurs:Algebraische Kurven (Osnabrück 2012)/Arbeitsblatt 27/latex
Kurs:Algebraische Kurven (Osnabrück 2012)/Arbeitsblatt 27/latex


\setcounter{section}{27}


  \zwischenueberschrift{Aufwärmaufgaben}


\inputaufgabe

{}

{


Definiere eine
\definitionsverweis {Äquivalenzrelation}{}{}
auf der Menge

\mathl{{ {\mathbb A}_{  K  }^{ n+1  } } \setminus \{0\}}{} derart, dass der Quotient unter der Äquivalenzrelation der projektive $n$-dimensionale Raum ist.


}

{} {}


\inputaufgabe

{}

{


Man definiere den Begriff \stichwort {projektiv-linearer Unterraum} {} eines projektiven Raumes ${\mathbb P}^{ n }_{ K }$.


}

{} {}


\inputaufgabe

{}

{


Es sei


\mavergleichskette

{\vergleichskette

{ {\mathfrak a}
}

{ \subseteq }{ K[X_1 , \ldots ,X_n]
}

{  }{
}

{    }{
}

{   }{
}

}
{}{}{}
ein
\definitionsverweis {Ideal}{}{.}
Zeige, dass ${\mathfrak a}$ genau dann ein
\definitionsverweis {homogenes Ideal}{}{}
ist, wenn es von homogenen Elementen erzeugt wird.


}

{} {}


\inputaufgabe

{}

{


Es sei


\mavergleichskette

{\vergleichskette

{K
}

{ = }{ {\mathbb F}_{  q   }
}

{  }{
}

{    }{
}

{   }{
}

}
{}{}{}
ein
\definitionsverweis {endlicher Körper}{}{}
mit $q$ Elementen. Berechne auf zwei verschiedene Arten, wie viele Elemente der
\definitionsverweis {projektive Raum}{}{}
${\mathbb P}^{ n }_{ K}$ besitzt.


}

{} {}


Die folgenden drei Aufgaben besprechen die Zariski-Topologie auf den projektiven Räumen.


\inputaufgabe

{}

{


Zeige, dass die
\definitionsverweis {Zariski-Topologie}{}{}
auf dem
\definitionsverweis {projektiven Raum}{}{}
wirklich eine
\definitionsverweis {Topologie}{}{}
ist.


}

{} {}


\inputaufgabe

{}

{


Es sei $K$ ein unendlicher
\definitionsverweis {Körper}{}{}
und ${\mathbb P}^{ n }_{ K}$ der
\definitionsverweis {projektive Raum}{}{.}
Charakterisiere die
\definitionsverweis {homogenen Ideale}{}{}
${\mathfrak a}$, für die


\mavergleichskette

{\vergleichskette

{ D_+({\mathfrak a})
}

{ = }{ \emptyset
}

{  }{
}

{    }{
}

{   }{
}

}
{}{}{}
ist.


}

{} {}


\inputaufgabe

{}

{


Es sei $K$ ein unendlicher Körper. Zeige, dass der projektive Raum ${\mathbb P}^{ n }_{ K}$
\definitionsverweis {irreduzibel}{}{}
ist.


}

{} {}


  \zwischenueberschrift{Aufgaben zum Abgeben}


\inputaufgabe

{}

{


Zeige, dass zwei verschiedene Punkte $P$ und $Q$ in der projektiven Ebene eindeutig eine projektive Gerade definieren, auf der beide Punkte liegen. Wie berechnet man die Geradengleichung aus den Koordinaten der Punkte?


}

{Bestimme die homogene Geradengleichung für die beiden Punkte

\mathl{(2,3,7)}{} und

\mathl{(1,5,-2)}{.}} {}


\inputaufgabe

{}

{


Es sei ${\mathbb P}^{ n }_{ K }$ ein
\definitionsverweis {projektiver Raum}{}{}
der Dimension $n$ und es seien


\mavergleichskette

{\vergleichskette

{ X,Y
}

{ \subseteq }{ {\mathbb P}^{ n }_{ K }
}

{  }{
}

{    }{
}

{   }{
}

}
{}{}{}
projektiv-lineare Unterräume der Dimension $r$ und $s$. Es sei


\mavergleichskette

{\vergleichskette

{ r+s
}

{ \geq }{ n
}

{  }{
}

{    }{
}

{   }{
}

}
{}{}{.}
Zeige, dass dann


\mavergleichskette

{\vergleichskette

{ X \cap Y
}

{ \neq }{ \emptyset
}

{  }{
}

{    }{
}

{   }{
}

}
{}{}{}
ist.


}

{} {}


\inputaufgabe

{}

{


Es sei $K$ ein unendlicher
\definitionsverweis {Körper}{}{}
und sei

\mathl{P_1 , \ldots , P_m}{} eine endliche Ansammlung von Punkten in einem
\definitionsverweis {projektiven Raum}{}{}
${\mathbb P}^{ n }_{ K }$. Zeige: Dann gibt es eine homogene Linearform


\mavergleichskette

{\vergleichskette

{ L
}

{ \in }{ K[X_0 , \ldots , X_n]
}

{  }{
}

{    }{
}

{   }{
}

}
{}{}{}
derart, dass all diese Punkte auf der durch $L$ definierten offenen Teilmenge $D_+(L)$ liegen.


}

{} {}


Die nächste Aufgabe benötigt noch die folgende Definition:


Für ein
\definitionsverweis {homogenes Ideal}{}{}
$I$ in


\mavergleichskette

{\vergleichskette

{ R
}

{ = }{ A[X_0 , \ldots , X_n]
}

{  }{
}

{    }{
}

{   }{
}

}
{}{}{}
mit der
\definitionsverweis {Standardgraduierung}{}{}
definiert man die \definitionswort {Sättigung}{}
\zusatzklammer {oder \definitionswort {Saturierung}{}} {} {}
von $I$ als


\mathdisp {\left\{ r \in R  \mid   \text{es existiert ein } n \text{ mit } r \cdot (R_+)^n \subseteq I         \right\}} { . }


Dabei ist $R_+$ das
\definitionsverweis {irrelevante Ideal}{}{}


\mavergleichskette

{\vergleichskette

{  \bigoplus_{d \geq 1}R_d
}

{ = }{ \left(  X_0 , \ldots , X_n  \right)
}

{  }{
}

{    }{
}

{   }{
}

}
{}{}{.}


\inputaufgabe

{}

{


Es sei $A$ ein
\definitionsverweis {kommutativer Ring}{}{}
und


\mavergleichskette

{\vergleichskette

{ R
}

{ = }{ A[ X_0 , \ldots ,  X_n]
}

{  }{
}

{    }{
}

{   }{
}

}
{}{}{}
der
\definitionsverweis {Polynomring}{}{}
mit der
\definitionsverweis {Standardgraduierung}{}{.}
Zeige, dass die
\definitionsverweis {Sättigung}{}{}
eines
\definitionsverweis {homogenen Ideals}{}{}
$I$ wieder ein homogenes Ideal ist.


}

{} {}


<< | Kurs:Algebraische Kurven (Osnabrück 2012) | >>


PDF-Version dieses Arbeitsblattes


 Zur Vorlesung
(PDF)
